% 328_Kopplungsregime_Resonanz_En_ch.tex
% Doc. 328, FFGFT corpus, August 2026 — chapter file
% Coupling Regimes, Resonance, and Synchronisation

% --- Local macros ---
\providecommand{\K}{}\renewcommand{\K}{\textbf{[K]}}
\providecommand{\B}{}\renewcommand{\B}{\textbf{[B]}}
\renewcommand{\S}{\textbf{[S]}}
\providecommand{\SET}{}\renewcommand{\SET}{\textbf{[AXIOM]}}
\providecommand{\E}{}\renewcommand{\E}{\textbf{[E]}}
\providecommand{\kk}{k_{\mathrm{crit}}}
\providecommand{\kmix}{\kappa_{\mathrm{mix}}}
\providecommand{\Keff}{K_{\mathrm{eff}}}
\providecommand{\xipar}{\xi}

\definecolor{ffgftblue}{RGB}{30,90,160}
\definecolor{proofgreen}{RGB}{0,110,50}
\definecolor{warnviolet}{RGB}{100,0,130}
\definecolor{etabliert}{RGB}{0,80,120}
\tcbset{
  base/.style={breakable,enhanced,boxrule=0.6pt,arc=3pt,
    fonttitle=\bfseries\small,coltitle=white,top=4pt,bottom=4pt,left=6pt,right=6pt},
  ffgftbox/.style={base,colback=blue!5,colframe=ffgftblue,colbacktitle=ffgftblue},
  proofbox/.style={base,colback=green!5,colframe=proofgreen,colbacktitle=proofgreen},
  warnbox/.style={base,colback=violet!6,colframe=violet!60!black,colbacktitle=violet!60!black},
  etabliertbox/.style={base,colback=cyan!5,colframe=etabliert,colbacktitle=etabliert}
}

% ============================================================
\section*{Abstract}

Communications engineering has a quantitative three-way classification of coupling
regimes (undercritical, critical, overcritical) relative to the damping of the
coupled subsystems. This document shows that this classification is an overlooked
structural principle: the same mathematical structure recurs in quantum optics,
particle physics, and synchronisation theory. The connection to FFGFT is established
through ten verification scripts with explicit symbolic and numerical results.

\medskip
\noindent\textbf{Main results:}
\begin{itemize}
  \item \B\ Dictionary two-circuit $\leftrightarrow$ 2$\times$2 mixing (Hermitian +
    non-Hermitian, exceptional point $g_{\mathrm{EP}} = |\gamma_1-\gamma_2|/4$)
  \item \B\ Exact duality $D^2+V^2=1$ for all pure pointers; measurement transition
    as one-parameter continuum $\kappa = d/(2\sigma)$
  \item \K\ $\kmix(\theta_W)$ derived from $\xi$ by composition (Doc.\,323);
    screening negative (no simple constant)
  \item \B\ GST texture $H_{12}=\sqrt{m_1 m_2}$ proven symbolically (Doc.\,041 formulas)
  \item \K\ $\Keff = 2\pi\xi \approx 8.4\times10^{-4}$ from Doc.\,007/008;
    $\varphi$ quantitatively outside all Arnold tongues
  \item Four check-points identified in Doc.\,041 (author decision required)
\end{itemize}

% ============================================================
\section{Motivation and Thesis}

Theoretical physics treats couplings primarily via the magnitude of coupling constants.
Communications engineering, by contrast, has an older and structurally richer
classification: couplings are judged not absolutely but \emph{relative to the damping
of the coupled subsystems}. This yields a three-way split — undercritical, critical,
overcritical — distinguishing not merely quantitative shifts but \textbf{qualitative
changes of spectral structure}.

\begin{tcolorbox}[ffgftbox,title={Thesis}]
This regime classification is an overlooked ordering principle. It recurs in quantum
optics (strong coupling), particle physics (mixing, avoided crossing), and
synchronisation theory (Kuramoto threshold, Arnold tongues) under different names.
For FFGFT, in which particles are resonance states, it provides the natural language
for when one state becomes two — and which frequency ratios lock or remain stably
incommensurate.
\end{tcolorbox}

% ============================================================
\section{The Two-Circuit Formalism \E}

Two resonators with eigenfrequencies $\omega_1$, $\omega_2$, quality factors $Q_1$, $Q_2$
and coupling coefficient $k$. Reference quantity: \textbf{critical coupling}
\begin{equation}
  \kk = \frac{1}{\sqrt{Q_1 \cdot Q_2}}.
\end{equation}
Normalised coupling factor $\kappa = k/\kk = k\cdot Q$ (symmetric case):

\begin{itemize}
  \item \textbf{Undercritical} ($\kappa < 1$): single peak at $\omega_0$; incomplete energy transfer.
  \item \textbf{Critical} ($\kappa = 1$): maximally flat pass-band (Butterworth); optimal power transfer.
  \item \textbf{Overcritical} ($\kappa > 1$): two peaks at
    $\omega_\pm \approx \omega_0\!\left(1 \pm \tfrac{1}{2}\sqrt{k^2-\kk^2}\right)$.
\end{itemize}

\begin{tcolorbox}[etabliertbox,title={Two structural statements \E}]
\textbf{(2.1)} Coupling shifts frequencies before splitting them (frequency pulling);
the ``bare'' eigenfrequency is in principle unobservable.\\[2pt]
\textbf{(2.2)} The threshold is a ratio: whether a system splits depends on $k$
relative to the loss rates, not on $k$ alone.
\end{tcolorbox}

\noindent Numerically verified by \texttt{pruef\_328\_2\_bandfilter.py}:
peak formula to $<10^{-5}$ relative, Butterworth curvature exactly zero at $\kappa=1$.

% ============================================================
\section{Recurrence in Quantum Physics \E}

The correspondence is mathematically exact, not merely analogical:

\begin{center}
\begin{tabular}{lll}
\toprule
\textbf{Communications eng.} & \textbf{Quantum physics} & \textbf{Condition}\\
\midrule
Overcrit. peak splitting & Rabi splitting, avoided crossing & $g > (\kappa_c+\gamma)/2$\\
Critical coupling & Critical coupling (resonator QED) & coupling $=$ loss rate\\
Frequency pulling & Mass/energy renormalisation, Lamb shift & coupling to continuum\\
Envelope interference & $\rho$–$\omega$ mixing & two poles, shared channel\\
\bottomrule
\end{tabular}
\end{center}

% ============================================================
\section{Measurement as Coupling — the $\kappa$-Continuum \B}

Every measurement couples the readout system to the object. Complete decoupling
$\equiv$ zero information flow.

\subsection{Symbolic Proof (von Neumann pointer model)}

System $a|+\rangle + b|-\rangle$, Gaussian pointer of width $\sigma$, shift $\pm d/2$;
coupling parameter $\kappa \equiv d/(2\sigma)$. Exact results:

\begin{tcolorbox}[proofbox,title={Exact duality — \B\ (Scripts 5 and 9A)}]
\begin{align}
  V(\kappa) &= \langle\phi_+|\phi_-\rangle = e^{-\kappa^2/2}
    \quad\text{(back-action / residual coherence)},\\
  D(\kappa) &= \sqrt{1 - e^{-\kappa^2}}
    \quad\text{(Helstrom distinguishability)},\\
  \boldsymbol{D^2 + V^2} &= \boldsymbol{1}
    \quad\text{for \emph{every} pure pointer.}
\end{align}
Limits: $\kappa\to0$: $(V,D)=(1,0)$; $\kappa\to\infty$: $(V,D)=(0,1)$.\\
Weak regime: $D\approx\kappa$ (linear), $1-V\approx\kappa^2/2$ (quadratic).
\end{tcolorbox}

\begin{tcolorbox}[warnbox,title={Not closable — \SET}]
``Every measurement apparatus reduces to a pointer model'' is a physical model
assumption, not provable, only falsifiable. Declared as \SET\ (model framework).
\end{tcolorbox}

All quantities $V,D,\kappa$ reside on the state-transition level of the A271
three-level separation (hardware / state transition / description).

% ============================================================
\section{Synchronisation and Locking \E}

\subsection{Adler Equation}
$\dot{\phi} = \Delta\omega - \varepsilon\sin\phi$. Lock-in precisely for $|\Delta\omega|\le\varepsilon$;
lock range $\propto$ coupling $\varepsilon$. Verified numerically
(\texttt{pruef\_328\_3\_einrasten.py}, deviation $<0.1\,\%$).

\subsection{Kuramoto Model}
Below the critical coupling $K_c$ the ensemble remains incoherent ($r=0$);
above $K_c$ the order parameter $r\propto\sqrt{K-K_c}$ grows continuously —
a second-order phase transition. This is \textbf{not} a three-way split in the
sense of Sec.\,2: the Kuramoto transition separates two structural regimes
(disorder / growing order), not three. The closer parallel to the band-filter
three-way classification is the Adler model (Sec.\,5.1): there the lock-range
threshold is exactly the coupling amplitude, sharing the same threshold structure
as the band-filter critical coupling.

\subsection{Noise Minimum}
The locked state exports entropy and is the least noisy available state —
the foundation of every frequency synthesis (PLL).

\subsection{Thermodynamic Clarification}
Phase locking occurs exclusively in open, driven, dissipative systems. ``Natural
systems tend toward an optimum, not toward uniform distribution'' holds for
nonequilibrium steady states; the second law remains intact.

% ============================================================
\section{Arnold Tongues, KAM, and the Golden Ratio}

\subsection{Qualitative stability \E}
Tongue widths $\Delta\Omega(p/q)\propto K^q$ (Script 9B: slope 1.98 for $q=2$,
2.69 for $q=3$). The most irrationally approximable ratio is $\varphi=[1;1,1,\ldots]$
(KAM). Connects to Doc.\,172/188/275 of the corpus.

\subsection{Quantitative FFGFT connection \K}
From Doc.\,008: $m(n)=m_0 e^{\xi n}$, coupling term $g\sim\xi\omega$, hence
$g/\omega=\xi$ ($n$-independent). Effective $K$ in the circle map:
\begin{equation}
  \Keff = 2\pi\xi \approx 8.4\times10^{-4}.
\end{equation}
The $\xi^2$ suppression (Doc.\,007: $\sigma_\nu\sim\xi_0^2 G_F$) corresponds to the
$K^2$ term in the tongue width. At $\Keff\ll1$:
$\Delta\Omega(1/2)\approx 5.6\times10^{-8}$.
$\varphi$ lies entirely outside every tongue. \K\ (Orbifold derivation outstanding.)

% ============================================================
\section{Mixing Angles as Coupling Grades — Largely Closed}

\subsection{Dictionary two-circuit $\leftrightarrow$ 2$\times$2 mixing \B}

Hermitian case (Script 1):
\begin{equation}
  \lambda_\pm = \tfrac{H_{11}+H_{22}}{2} \pm
  \sqrt{\!\left(\tfrac{H_{11}-H_{22}}{2}\right)^2+H_{12}^2},\quad
  \tan 2\theta = \frac{2H_{12}}{H_{11}-H_{22}}.
\end{equation}

Non-Hermitian extension (Script 8, exceptional point):
\begin{equation}
  g_{\mathrm{EP}} = \frac{|\gamma_1-\gamma_2|}{4}.
\end{equation}
For $g<g_{\mathrm{EP}}$: only widths split (undercritical);
for $g>g_{\mathrm{EP}}$: frequencies split (Rabi, overcritical).
\textbf{Band-filter $\kk$ and cavity-QED ``strong coupling'' are special cases of
this EP condition.} \B

\begin{center}
\small
\begin{tabular}{llll}
\toprule
\textbf{Angle} & $\theta$ [°] & $\sin^2\!\theta$ & $\kmix = \tan 2\theta$\\
\midrule
$\theta_W$ (Doc.\,323: $\sin^2\theta_W=0.2308$) & 28.71 & 0.2308 & 1.57\\
Cabibbo $\theta_C$ ($|V_{us}|=0.22452$, Doc.\,041) & 12.98 & 0.050 & 0.49\\
CKM $\theta_{13}$ ($|V_{ub}|=0.00365$) & 0.21 & $\approx0$ & 0.007\\
PMNS $\theta_{12}$ (Doc.\,041: 33.44°) & 33.44 & 0.304 & 2.34\\
PMNS $\theta_{23}$ (Doc.\,041: 49.2°) & 49.2 & 0.573 & $-6.77$\\
PMNS $\theta_{13}$ (Doc.\,041: 8.57°) & 8.57 & 0.022 & 0.31\\
\bottomrule
\end{tabular}
\end{center}

\subsection{$\xi$-Level: $\theta_W$ sector closed \K}

$\sin^2\theta_W=0.2308$ derived from $\xi$ in Doc.\,323 \K. Hence
$\kmix(\theta_W)=1.5652$ is $\xi$-derived by composition. Screening against 11
candidate constants: no hit (closest: $e/\!\sqrt{3}$, $2.7\times10^{-3}$).

\subsection{GST texture and $H_{12}$ \B}

Doc.-041 formulas of type $\theta=\arcsin\!\sqrt{m_1/m_2}$ are exactly the
Gatto–Sartori–Tonin texture (Script 7):
\begin{equation}
  M = \begin{pmatrix}0 & \sqrt{m_1 m_2}\\ \sqrt{m_1 m_2} & m_2-m_1\end{pmatrix},
  \quad \text{eigenvalues}\;\{-m_1,m_2\},\quad \tan\theta=\sqrt{m_1/m_2}.
\end{equation}
$H_{12}=\sqrt{m_1 m_2}$ (geometric mean of masses). \B

\subsection{Four check-points from Doc.\,041}

\begin{tcolorbox}[warnbox,title={Check-points for Doc.\,190 — author decision required}]
(i)~$\delta_{\rm CKM}=\arcsin(2\sqrt{2}\sqrt{\xi}/3)$ gives $0.011\,\mathrm{rad}$
    vs.\ $1.20\,\mathrm{rad}$; candidate: $\arcsin(2\sqrt{2}/3)=1.231\,\mathrm{rad}$.\\
(ii)~$\theta_{13}=\arcsin(\xi^{1/3})$ gives $2.93°$ vs.\ $8.57°$;
    candidate: $\arcsin((25\xi)^{1/3})=8.59°$.\\
(iii)~$|V_{us}|$ with Doc.\,006 masses: $5.4\,\%$ discrepancy.\\
(iv)~Doc.\,007 $\leftrightarrow$ Doc.\,041 tension already declared as
    \texttt{\string\begin\{warning\}} and \texttt{\string\begin\{speculation\}} in the source.
\end{tcolorbox}

% ============================================================
\section{Delimitation}

\begin{tcolorbox}[warnbox,title={Delimitation from resonance-as-vocabulary}]
Schepis (IPI Letters 3(2) 2025; ICIP-2026) uses resonance as connecting vocabulary
without computable quantities. This document takes the opposite direction: every
bridge is verified by a script and carries a definite status marker. Observers appear
only as read-out sub-circuits; no consciousness vocabulary.
\end{tcolorbox}

% ============================================================
\section{Closability Balance}

\begin{tabular}{lp{9.5cm}}
\toprule
\textbf{Bridge} & \textbf{Status}\\
\midrule
R-new-b & \B\ Hermitian + non-Hermitian (EP), complete\\
R-new-c & \K\ for $\theta_W$ (composition); \B\ GST texture; four check-points open\\
R-new-d & \B\ Duality $D^2+V^2=1$, all pure pointers\\
R-new-a (qual.) & \B\ $K^q$ hierarchy, $\varphi$ quantitatively outside tongues\\
R-new-a (quant.) & \K\ $\Keff=2\pi\xi$ from Doc.\,007/008; orbifold derivation outstanding\\
\SET & Reducibility measurement $\to$ pointer; in principle unprovable\\
\bottomrule
\end{tabular}

% ============================================================
\section{Verification Scripts}

\begin{tabular}{lll}
\toprule
\textbf{Script} & \textbf{Content} & \textbf{Result}\\
\midrule
pruef\_328\_1 & Dictionary, sympy & \B\\
pruef\_328\_2 & Band-filter regimes, numpy & \E\ confirmed\\
pruef\_328\_3 & Adler, Arnold tongues & \E\ confirmed\\
pruef\_328\_4 & $\kmix$ table & Finding\\
pruef\_328\_5 & Measurement $\kappa$, $D^2+V^2$ & \B\\
pruef\_328\_6 & $\xi$-composition, screening & \K\\
pruef\_328\_7 & GST texture, corpus check & \B, check-points\\
pruef\_328\_8 & EP, non-Hermitian & \B\\
pruef\_328\_9 & Pure pointers, $K^q$ & \B\\
pruef\_328\_10 & $\Keff=2\pi\xi$ & \K\\
\bottomrule
\end{tabular}
